\section{Discussion}\label{sec:discussion}

\paragraph{What a small model can and cannot do.} Finding the wearer's hand, the fingertip and whether a hand is present works at this size (hand P@0.5 of 92.6\%, 3.9\,px fingertip error, AUROC 0.998), although compact detectors draw the hand box more tightly. Deciding which object the finger refers to, and how far that object extends, is what large models do much better; compact detectors of the same compute pick the object slightly less often than SpatialNet but draw its extent much better, even with the same ImageNet pretraining (Section~\ref{sec:detectors}). EgoPoint is dominated by overlap: in about 70\% of images the fingertip lies inside the target box, so for most targets the fingertip tells \emph{where} the object is but not \emph{how large} it is (a mean-size box at the true fingertip reaches 11.4\% P@0.5). One possible reason is that people align the fingertip with the target from the eye~\cite{mayer2015modeling,plaumann2018towards}, so that a camera next to the eye sees the target at or just behind the fingertip; because 61\% of EgoPoint images are composed or generated (Section~\ref{sec:data}), this remains a hypothesis about real pointing. Large-gap targets, where the ray matters, are only 7\% of the data. That is too few to shape a learned ray prior, and it is where errors concentrate.

\paragraph{Design implications.} For pointing interfaces on glasses, our results suggest six practical rules.
\begin{enumerate}
\item \emph{Choose the model by what the interface needs.} For a tight target box, a compact detector trained on pointing data is more accurate at the same compute, even without COCO pretraining. For naming the target, a structured model did as well, and it also provides the hand, the fingertip and the direction for feedback or gating in the same pass (Section~\ref{sec:detectors}).
\item \emph{Train with distractor hands, whatever the model.} Pasting unlabelled hands into the training images makes SpatialNet and a generic detector equally stable when other hands are in view, at no cost on clean images (Section~\ref{sec:detectors}).
\item \emph{Spend extra compute on resolution first.} At similar extra compute, a larger input gained more than a larger backbone (Sections~\ref{sec:resolution} and~\ref{sec:ablations}).
\item \emph{Confirm large-gap selections}, for example by naming or highlighting the target before acting on it, because errors concentrate there; for these targets, a recogniser does about as well with the whole frame as with SpatialNet's crop (Sections~\ref{sec:detectors} and~\ref{sec:heads-results}).
\item \emph{Smooth over time before triggering.} Per-frame false alarms make single-frame triggering unreliable (Section~\ref{sec:heads-results}).
\item \emph{Keep frames on the device.} A later recogniser needs only the target crop (Section~\ref{sec:problem}), which limits computation, bandwidth and the bystander imagery that leaves the glasses. On our unit, continuous local inference also used less charge than streaming frames to the phone, although one-minute observations do not establish an energy advantage (Section~\ref{sec:ondevice}).
\end{enumerate}

\paragraph{Limitations.}
\begin{itemize}
\item \emph{Construct validity.} The gap categories are image-plane relations. The direction label runs from the fingertip to the target centre, not along the finger, so direction errors mix pointing geometry with the choice of target centre. The ray lives in the image plane, so objects at different depths along the same image ray are told apart only by appearance, and the eye--camera offset is ignored.
\item \emph{Internal validity.} The ablation, resolution and backbone variants are single runs whose checkpoints come from different epochs (54 to 120), and some interventions are bundled (A1, A2 and A4). Paired intervals capture test-sample variation, not training randomness, and the per-component reading of the ablations is exploratory. Early stopping ended the one-cycle schedule before it completed.
\item \emph{External validity.} All accuracy figures come from one dataset with mostly landscape, partly composed or generated images, 4.7\% of whose test images duplicate training or validation images, and with unknown demographic coverage (for example, skin tone and handedness; flipping covers left hands only geometrically). Frames are squashed to a square, and the glasses' camera is portrait. EgoPoint has no non-pointing hands and no pointing without a target, so the pointing and target gates are untested, and in deployment a non-pointing hand may be read as pointing. The model has no temporal context. Our distractors are synthetic and encode the prior that the wearer's arm enters from the bottom of the view; a raised arm, or a companion's hand near the bottom, can break it, and we did not test real frames with other people's hands. The check on frames from the glasses (Section~\ref{sec:realframes}) covers only the wearer's hand and is small.
\item \emph{Baselines and measurement.} We copied the published large-model results and did not re-run any multimodal LLM under our protocol; the compact detectors we did train draw better boxes at the same compute. COCO-pretrained detectors have seen the source photographs of at least 41\% of our test images, but COCO pretraining gave YOLO26n no extra gain on them; published models trained on large image collections may have seen them too, which we could not test. The remedies, the ablations and three of the four controlled detector runs used one training seed each. The recognition test is category-level, uses two zero-shot recognisers, and measures what a recogniser receives, not what a user experiences. Within EgoPoint itself, 230 test images share their source photograph with a training or validation image. The device measurements are single observations on a pre-production unit, and we did not measure the detectors on the glasses. Finally, we measure the model on a benchmark and on the device, not the interaction: how wearers respond to its errors, and which confirmation designs suit them, need a study with participants wearing the glasses.
\end{itemize}
